\zlabel{firstpagetocount}       % DO NOT REMOVE! Used for counting number of pages of main text

% Part Labels Explanation:
% Labels in the form \label{chapter:<part_name>} are used to categorize chapters into specific parts of the thesis structure.
% These labels are essential for calculating the ratio of text dedicated to different thesis sections (e.g., Introduction, Methodology, Results, Discussion, Summary).
% Allowed labels:
% - chapter:introduction -> Introduction (<10% of the text).
% - chapter:method       -> Methodology (<20% of the text).
% - chapter:results      -> Results (30–40% of the text).
% - chapter:discussion   -> Analysis, Discussion, and Conclusions (30–40% of the text).
% - chapter:summary      -> Summary (<0.5 pages).
% How to use: Add a \label{chapter:<part_name>} for each chapter to indicate its part. 
% These labels ensure consistency and allow automated tests to validate the thesis structure.

\chapter{Sissejuhatus}\label{chapter:introduction} % This command creates a numbered chapter titled "Sissejuhatus" (Estonian for "Introduction"). The \label{chapter:introduction} assigns a label to the chapter, allowing you to reference it later (e.g., \ref{chapter:introduction}).
\input{chapters/introduction} %  This command inserts the content from the file introduction.tex located in the chapters folder.

\chapter{Taust}\label{chapter:method}
\input{chapters/background.tex}

\chapter{Metoodika ja andmestik}\label{chapter:method}
\input{chapters/method.tex}

\chapter{Katsed ja hindamismeetodid}\label{chapter:discussion}
\input{chapters/experiments.tex}

\chapter{Andmeanalüüs ja demonstreerimine}\label{chapter:discussion}
\input{chapters/analysis.tex}

\chapter{Tulemused ja järeldused}\label{chapter:result}
\input{chapters/result.tex}

\chapter{Kokkuvõte}\label{chapter:summary} 
\input{chapters/summary.tex}

\zlabel{lastpagetocount}        % DO NOT REMOVE! Used for counting number of pages of main text
